\section{Two degree-six class certificates for the negative conductor-eight character}
\label{sec:negative-quadratic8-class-certificates}

Put $p=269$, $K=\mathbb Q(\sqrt{-3})$, and $t=\sqrt[3]p>0$.
For $d=2,538$ define
\[
 w_d=\sqrt{-d},\qquad F_d=\mathbb Q(w_d),\qquad
 L_d=F_d(t),\qquad M_d=KL_d.
 \tag{G150C.1}
\]
Let $\delta_2=\chi_{-8}$ be the primitive odd quadratic character
of conductor eight, and set $\delta_{538}=\chi_{-8}\tau_p$,
where $\tau_p=(\cdot/p)$.
Its fundamental conductor is $8p=2152$.
Characters are extended by zero at integers not coprime to their
primitive conductor.

\begin{theorem}[AC-G150-class: the two negative quadratic-eight class sectors]
\label{AC-G150-class}
For both fields in (G150C.1),
\[
                      \operatorname{Cl}(L_d)[269]=0.
 \tag{G150C.2}
\]
Consequently the nontrivial cubic sector of
$\operatorname{Cl}(M_d)[269]$ vanishes.
More precisely, two explicitly certified units generate a subgroup
of finite index $m_d$ in $\mathcal O_{L_d}^{\times}/\mu_{L_d}$,
with
\[
 h_{L_2}=42m_2,\qquad h_{L_{538}}=140m_{538},\qquad
 269\nmid m_2m_{538}.
 \tag{G150C.3}
\]
No assertion that either index is one, that either class number is
exactly the displayed selector, or that the full class group of
$M_d$ vanishes is made.
This is a class-sector theorem; the original-coefficient strict
group and the integral-index window require separate arguments.
\end{theorem}

\begin{proof}
\proofstep{1. Exact fields, local discriminants, and signatures.}
The fields $F_d$ have fundamental discriminants $-8$ and $-2152$,
respectively.  They are the fields of $\delta_d$.
The splitting field of $X^3-p$ has group $S_3$ and its unique
quadratic subfield is $K$, distinct from both $F_d$.
Thus
\[
 [L_d:\mathbb Q]=6,\quad
 \operatorname{Gal}(M_d/\mathbb Q)=S_3\times C_2,
 \qquad \operatorname{Gal}(M_d/K)=C_3\times C_2.
 \tag{G150C.4}
\]
Each $L_d$ is fixed by a reflection in the $S_3$ factor.
Its signature is $(0,3)$: the imaginary quadratic field has two
conjugate embeddings and the cubic field has three embeddings.
Every abelian subfield of $M_d$ lies in $KF_d$, the fixed field
of the commutator $C_3$; its intersection with $L_d$ is $F_d$.
Therefore any root of unity in $L_d$ lies in $F_d$.
Since neither quadratic field is $\mathbb Q(i)$ or $K$,
\[
                 |\mu_{L_d}|=2,\qquad
                 \operatorname{rank}\mathcal O_{L_d}^{\times}=2.
 \tag{G150C.5}
\]

The pure cubic field $T=\mathbb Q(t)$ has discriminant $-3p^2$.
Here is the local check needed below.  At $p$ the defining
polynomial is Eisenstein, and the ramification is tame of degree three.
At three the equation $z+3z^2+3z^3=-30$ has a Hensel solution,
so $-(1+3z)$ is a cube root of $p$ in $\mathbb Q_3$.
The other factor is $\mathbb Q_3(\sqrt{-3})$; the local
discriminant exponent is one.  At two $T$ is unramified.
The polynomial discriminant $-27p^2$ excludes every other
possible ramified prime.  These local exponents therefore give
the stated field discriminant.

At two both $F_d$ have discriminant exponent three, and $T$
is unramified, so finite-etale base change gives exponent nine in $L_d$.
At three $F_d$ is unramified and the cubic exponent is one,
so the exponent is two.
At $p$, $(-2/p)=-1$: $F_{2,p}$ is the unramified quadratic
extension, and $L_{2,p}$ has $e=3,f=2$.
Its tame discriminant exponent is $f(e-1)=4$.
For $d=538$, the quadratic field is ramified of degree two;
the compositum with the tame cubic has $e=6,f=1$ and exponent five.
Consequently
\[
 \begin{aligned}
 D_2:=|D(L_2)|&=2^9 3^2p^4=24128014791168,\\
 D_{538}:=|D(L_{538})|&=2^9 3^2p^5=6490435978824192.
 \end{aligned}
 \tag{G150C.6}
\]
The completion of $K$ at $p$ is unramified quadratic.
Thus $M_2$ has two $p$-places, each of absolute local degree six,
with $e=3,f=2$; $M_{538}$ has one $p$-place, with $e=6,f=2$.
These different local situations are retained for subsequent strict arguments.

\proofstep{2. Four exact units, with no complete-basis assumption.}
Use the ordered basis $t^a w_d^b$, $0\le a<3$, $0\le b<2$,
with coordinate index $2a+b$.
The first certified units are
\[
 U_{2,1}=-270+739t-108t^2,\qquad U_{538,1}=-U_{2,1}.
 \tag{G150C.7}
\]
The second units have numerator vectors
\[
 \begin{split}
 v_2=(&-3665130267800299960477,-2767560365829808262340,\\
      &-241195089010706662698,562049263035469722408,\\
      &125317590322106716938,-20653277382608917902),
 \end{split}
 \tag{G150C.8}
\]
and
\[
 \begin{split}
 v_{538}=(&23480405011964765599140732341,\\
          &1974469358703520803947161356,\\
          &4325362166431043692713434890,\\
          &-288742341066790749679483080,\\
          &-1233516577579131403410780250,\\
          &-2652837967636451338514550).
 \end{split}
 \tag{G150C.9}
\]
Their common denominators are respectively one and $807=3p$.
In particular the latter denominator must be retained in any
later local expansion.
The full exact certificates include their monic integral
degree-six relations with constant coefficient one.
The input hashes are \texttt{01228a61...} and \texttt{ea592df7...};
the evidence manifest records all complete hashes and coefficients.

For clarity, the verification does not use a provisional integral
basis or class group.  The six-by-six rational multiplication
matrices $T,W$ satisfy
\[
              T^3=pI,\quad W^2=-dI,\quad TW=WT.
 \tag{G150C.10}
\]
The raw generator is $w_d-t$.
Its exact characteristic polynomial is the saved resultant,
which is irreducible over $\mathbb Q$ and has no real root.
Both inverse rational maps for $t$ and $w_d$ are verified.
Every candidate unit is evaluated at $W-T$; its first column
is exactly the displayed tensor vector divided by its denominator.
Its characteristic polynomial is the supplied monic integral
polynomial with constant one, and a separate Horner evaluation
of that polynomial at the matrix is zero.
It follows that the unit and its inverse are algebraic integers.
All four objects are therefore actual global units independently
of the proposal-only extraction.

\proofstep{3. Residue logarithms and the original unit index.}
Take $\ell=206593$, so $\ell-1=768p$.
Complete trial division verifies primality; the power tests for
the prime factors of $\ell-1$ verify that $g=10$ is primitive.
Put $h=g^{768}=49873\pmod\ell$, of order $p$.
The three roots of $X^3-p$ and the two roots of $X^2+d$ are
simple modulo $\ell$, and give six degree-one $\ell$-adic
embeddings by Hensel lifting.
All denominators in Step2 are invertible modulo $\ell$.
For every embedding define $c_j\in\mathbb F_p$ by
$\overline{U}_{d,j}^{768}=h^{c_j}$.
The six residue rows are saved; the first two give
\[
 C_2=\begin{pmatrix}85&173\\85&96\end{pmatrix},\qquad
 C_{538}=\begin{pmatrix}85&109\\85&160\end{pmatrix}.
 \tag{G150C.11}
\]
Their determinants are $-6545\equiv180$ and $4335\equiv31$
modulo $p$, respectively.
These are exact finite logarithms, obtained by enumerating the
$269$ powers of $h$; no real or $p$-adic logarithm approximation occurs.

The residue maps kill $p$th powers and the roots of unity.
Thus the two units are independent modulo $p$ in
$\mathcal O_{L_d}^{\times}/\mu_{L_d}$.
They are independent over $\mathbb Z$ as well: a primitive
integer relation would remain nonzero modulo $p$.
The unit theorem in (G150C.5) makes their subgroup full rank,
and the same mod-$p$ invertibility proves
\[
 m_d=[\mathcal O_{L_d}^{\times}/\mu_{L_d}:
             \langle U_{d,1},U_{d,2}\rangle]<\infty,
                    \qquad p\nmid m_d.
 \tag{G150C.12}
\]

\proofstep{4. Weighted regulators and a class-number denominator bound.}
Let $\rho=p^{1/3}$ and $\omega=e^{2\pi i/3}$.
The three complex-place representatives are
$(t,w_d)=(\rho\omega^j,i\sqrt d)$ for $j=0,1,2$.
Form the three-by-two matrix $2\log|\sigma_j(U_{d,k})|$;
delete its last row and take the absolute determinant $R_{U,d}$.
This is the projected weighted-log regulator convention, so
\[
                         R_{U,d}=m_dR_{L_d}.
 \tag{G150C.13}
\]
There is no Euclidean hyperplane-covolume factor.
Evaluation of the exact tensor arrays with $1500$-bit outward
balls gives positive determinants with
\[
 3678<R_{U,2}<3679,\qquad 4353<R_{U,538}<4354.
 \tag{G150C.14}
\]
Every weighted column sum contains zero, as required by the norm.
The original high-precision regulator balls are retained in the
later $256$-bit analytic calculations.

For $n=6,r_2=3$ put
\[
 U_d=\left\lceil\frac{6!4^3(\lfloor\sqrt{D_d}\rfloor+1)}
                            {6^6 3^3}\right\rceil,\qquad
 H_d=4^6\left\lceil\sqrt[3]{U_d^4}\right\rceil.
 \tag{G150C.15}
\]
Minkowski and $\pi>3$ give one integral representative of each
ideal class of norm at most $U_d$.
Counting these ideals by their $4/3$-Dirichlet series gives
\[
 h_{L_d}\le U_d^{4/3}\zeta_{L_d}(4/3)
 \le U_d^{4/3}\zeta(4/3)^6\le4^6 U_d^{4/3}\le H_d,
 \tag{G150C.16}
\]
where $\zeta(4/3)\le1+\int_1^\infty x^{-4/3}dx=4$.
The Euler-factor comparison uses at most six prime ideals over
each rational prime, with residue degrees at least one.
Exact integer arithmetic gives
\[
 \begin{array}{c|r|r}
 d&U_d&H_d\\\hline
 2&179681&41530085376\\
 538&2946987&1730603491328.
 \end{array}
 \tag{G150C.17}
\]
If $r=h_{L_d}/m_d>1$ and $|r-c|<1/H_d$ for an integer $c$,
then $m_d<h_{L_d}\le H_d$.
Unless $h_{L_d}=cm_d$, the nonzero integer $h_{L_d}-cm_d$
would force $|r-c|\ge1/m_d>1/H_d$.
This isolates the rational ratio without assuming a complete unit basis.

\proofstep{5. Artin factors and the exact odd quadratic constants.}
Let $\varrho$ be the standard representation of the $S_3$
splitting field of $X^3-p$.
The permutation representation for $L_d$ is
$(1\oplus\varrho)\otimes(1\oplus\delta_d)$.
Artin formalism, including every inertia-invariant Euler factor,
therefore gives
\[
 Z_d:=\operatorname*{res}_{s=1}\zeta_{L_d}(s)
       =L(1,\delta_d)L(1,\varrho)L(1,\varrho\delta_d).
 \tag{G150C.18}
\]
For $f=8$ or $2152$ define
$B_{1,\delta}=f^{-1}\sum_{a=1}^{f-1}a\delta(a)$.
The exact finite sums are
\[
 B_{1,\delta_2}=-1,\qquad B_{1,\delta_{538}}=-10.
 \tag{G150C.19}
\]
Here $\chi_{-8}(a)$ is $+1$ for $a\equiv1,3\pmod8$,
$-1$ for $a\equiv5,7\pmod8$, and zero on even integers.
For $d=538$ it is multiplied by the exact Legendre symbol
$a^{134}\pmod p$.

The primitive finite Gauss sum is purely imaginary with squared
absolute value $f$.  The finite sine sum
$\sum_a\delta(a)\sin(2\pi a/f)$ is strictly positive in the
outward-ball calculation, fixing the sign as $i\sqrt f$.
The finite Fourier identity then gives
$\sum_a\delta(a)\sin(2\pi na/f)=\sqrt f\,\delta(n)$.
For $0<x<1$ the Abel limit of the sine series is
$\sum_{n\ge1}\sin(2\pi nx)/n=\pi(1/2-x)$.
Multiply by $\delta(a)$, sum over the finite set of residues,
and take the Abel limit of the convergent character Dirichlet series.
Since $\sum_a\delta(a)=0$, this proves
\[
 L(1,\delta_d)=-\frac{\pi B_{1,\delta_d}}{\sqrt f}
 =\begin{cases}\pi/\sqrt8,&d=2,\\10\pi/\sqrt{2152},&d=538.
 \end{cases}
 \tag{G150C.20}
\]
No quadratic class number or complete unit group is used here.

\proofstep{6. Exact Hecke coefficients, including the wild conductor at two.}
Write $L(s,\varrho)=\sum a_n n^{-s}$.
The coefficients are multiplicative, $a_1=1$,
$a_{3^v}=1$ and $a_{p^v}=0$ for $v\ge1$.
For primes $\ell\ne3,p$,
\[
 a_\ell=\#\{x\bmod\ell:x^3=p\}-1,\qquad
 a_{\ell^{v+1}}=a_\ell a_{\ell^v}-(-3/\ell)a_{\ell^{v-1}}.
 \tag{G150C.21}
\]
For $\ell\equiv2\pmod3$, the prime-power coefficients are
one in even degree and zero in odd degree.
For $\ell\equiv1\pmod3$ the exact cubic-residue power test
decides between $a_\ell=2$ and $-1$.
The local algebra at three in Step1 gives the root one,
while the cubic inertia at $p$ has no invariant vector.
The twisted coefficients are exactly $a_n\delta_d(n)$.
At two the twisting character is ramified and $\varrho$ is
unramified, so its inertia invariants vanish.
At $p$, its cubic characters remain nontrivial after the
possible quadratic twist; again there are no invariants.
Thus extension by zero gives the actual bad-prime factors.

Over $K$, write $\varrho=\operatorname{Ind}_K^{\mathbb Q}\nu$.
The finite cubic Kummer character $\nu$ has conductor $(p)$:
it is tame at $p$ and trivial locally at three because $p$
is a cube in $\mathbb Q_3$.
The twist is the norm pullback of $\delta_d$.
At two $K_2/\mathbb Q_2$ is unramified quadratic, and
\[
 N(1+4a)\equiv1+4\operatorname{Tr}_{\mathbb F_4/\mathbb F_2}
                                      (\bar a)\pmod8.
 \tag{G150C.22}
\]
The trace is onto and $\chi_{-8}(5)=-1$; hence the pullback
has exact conductor exponent three, not a squarefree conductor.
At $p$ the residue norm is onto, and multiplication of the
cubic character by a nontrivial quadratic character gives
order six, still tame with exponent one.
For $d=2$ the quadratic factor is unramified there.
Both twists remain unramified at three.
The primitive Hecke moduli and their rational conductors are therefore
\[
                    f=p,8p,\qquad Q=3f^2=217083,13893312.
 \tag{G150C.23}
\]

\proofstep{7. Primitive theta summation, phase and explicit tails.}
The imaginary quadratic Minkowski bound for $K$ is
$2\sqrt3/\pi<2$, so every ideal class has a representative of
integral norm one and $h(K)=1$.
Represent either primitive Hecke character by a multiplicative
character $\nu_f$ on $\mathcal O_K/f\mathcal O_K$, zero on
nonunits and trivial on the six global units.
Put $a=2\pi/(\sqrt3 f)$ and
\[
 \Theta(t)=\frac16\sum_{\alpha\in\mathcal O_K}
       \nu_f(\alpha)e^{-a|\alpha|^2t}
            =\sum_{n\ge1}b_ne^{-ant}.
 \tag{G150C.24}
\]
The second equality sums the six generators of each ideal.

The additive pairing
\[
 (\alpha,\beta)\longmapsto
 \exp\left(2\pi i\operatorname{Tr}_{K/\mathbb Q}
                    \frac{\alpha\bar\beta}{f\sqrt{-3}}\right)
 \tag{G150C.25}
\]
is perfect by the inverse different.
At unit frequency the finite Fourier sum is
$G\bar\nu_f(\bar\beta)$ by substitution.
At a nonunit frequency at a factor $\mathfrak q^a$ of the
modulus, choose a unit in $1+\mathfrak q^{a-1}$ on which the
primitive character is nontrivial, and one at the other factors;
for $a=1$ choose a residue unit instead.
It fixes the additive character but multiplies the Fourier sum
by a nontrivial scalar, so the sum is zero.
Parseval on the $f^2$ residue classes gives $|G|=f$.
This argument includes the factor $2^3$.

The Euclidean area of $f\mathcal O_K$ is $f^2\sqrt3/2$, with
dual lattice $2/(f\sqrt{-3})\mathcal O_K$.
Writing $b=2/\sqrt{-3}$ gives
$\operatorname{Re}(\alpha\overline{b\beta})
=\operatorname{Tr}(-\alpha\bar\beta/\sqrt{-3})$.
Thus the factor two in the dual cancels the one-half in the
real part; it is not presumed invertible modulo the even modulus.
Poisson summation for the Gaussian has scalar $1/(ft)$ and
dual exponent $a|\beta|^2/t$.
Conjugating the summation variable turns
$\bar\nu_f(\bar\beta)$ into $\bar\nu_f(\beta)$ without
changing the norm, and gives
\[
 \Theta(t)=\varepsilon t^{-1}\overline{\Theta(1/t)},
                        \qquad |\varepsilon|=1.
 \tag{G150C.26}
\]
Both constant terms vanish.
Mellin integration, initially for $\operatorname{Re}s>1$ and
then by exponential decay at both ends, gives at $s=1$
\[
 L(1)=\sum_{n\ge1}\frac{b_ne^{-an}}n
       +\varepsilon a\sum_{n\ge1}\bar b_n E_1(an),\qquad
 E_1(x)=\int_1^\infty e^{-xt}\frac{dt}{t}.
 \tag{G150C.27}
\]
All coefficient sequences here are real, and the computed theta
ball at one excludes zero.  Equation (G150C.26) at one therefore
forces the actual phase to be $\varepsilon=1$; no guessed
root-number convention is used.

The recurrence gives $|b_n|\le d(n)\le n$.
With $z=e^{-a}$ and $A_N=e^{-a(N+1)}$, the omitted tails of
the first sum, the second sum before its factor $a$, and theta
are bounded respectively by
\[
 \frac{A_N}{1-z},\qquad \frac{A_N}{a(1-z)},\qquad
                  \frac{A_N((N+1)-Nz)}{(1-z)^2}.
 \tag{G150C.28}
\]
The second estimate uses $E_1(x)\le e^{-x}/x$.
The script adds these full radii using $256$-bit outward balls
and cutoffs $12000$ and $80000$ at the two conductors.
It saves all weighted-log rows, finite Gauss sign balls,
theta values, tails, $L$-values and resulting intervals.

\proofstep{8. Rational isolation and precisely the reflected class sector.}
The analytic class-number formula for signature $(0,3)$ and
two roots of unity is
\[
 h_{L_d}R_{L_d}=\frac{2\sqrt{D_d}}{(2\pi)^3}Z_d.
 \tag{G150C.29}
\]
The certified inclusion calculations give, for $c_2=42$ and
$c_{538}=140$,
\[
 \frac{h_{L_d}R_{L_d}}{R_{U,d}}=\frac{h_{L_d}}{m_d}>1,
 \qquad \left|\frac{h_{L_d}}{m_d}-c_d\right|<6\cdot10^{-54},
 \qquad 1/H_d>5.7\cdot10^{-13}.
 \tag{G150C.30}
\]
The actual comparisons use the full balls, not these rounded
bounds.  The two integers enter only as final selectors, after
all field, unit, regulator, conductor and error calculations.
Step4 proves (G150C.3) exactly, and Step3 gives $p\nmid m_d$.
Since neither $42$ nor $140$ is divisible by $269$, this proves
the class assertion (G150C.2).

Let $s$ be the reflection fixing $L_d$ inside $M_d$.
Norm and extension on ideal classes have composites two and
$1+s$.  The norm of the $p$-primary class group is zero by
(G150C.2), so $1+s=0$ on it.
Every conjugate reflection therefore acts by $-1$ as well.
Their products generate the cubic rotation subgroup, which acts
trivially.  Thus exactly its nontrivial character sectors vanish;
the trivial cubic sector of $\operatorname{Cl}(M_d)$ is not asserted zero.
\end{proof}

\paragraph{Actual-character and original-coefficient boundary.}
The old finite labels under consideration are
$\chi^{-j}\chi_{-8}$, $j=1,2$, already contained in the
two-power cyclotomic part of the actual host.
At effective Tate residue $267$ the residual Kummer quotient is
$\bar\chi^{-j}\bar\chi_{-8}$, split by $M_2$;
at residue $133$ it is
$\bar\chi^{-j}\bar\chi_{-8}\bar\tau_p$, split by $M_{538}$.
The extra $\tau_p$ belongs only to residual splitting.
It is not inserted into the old finite host-character list.
The corresponding strict proof must still establish the local
unit injection, using the respective two-place and one-place
completions and every rational denominator, and then lift on
the originally given coefficient DVR.
This class theorem changes no coefficient ring, assumes no
absolute ramification restriction on it, and supplies no local
unit-injectivity or integral-index assertion by itself.

\paragraph{Primary inputs and computational scope.}
The analytic class-number formula and projected weighted-log
normalization are the retained Sutherland MIT 18.785 Lecture19
(2019), pp.1 and4,
\url{https://math.mit.edu/classes/18.785/2019fa/LectureNotes19.pdf}.
The Minkowski ideal bound, Dirichlet unit theorem, Artin formalism
and ordinary Poisson summation retain their previously audited scopes.
The odd quadratic formula is derived in Step5 from finite Fourier
summation and the Abel sine series, including an actual finite
Gauss-sign check; no quadratic class-number input is used.
The primitive prime-power theta argument and denominator isolation
are proved here with the new exact conductors and degree six.
All finite identities use exact integer, rational or modular arithmetic.
The inclusion computations use python-flint0.9.0 and its documented
outward Arb operations, including the order-one exponential integral;
the implementation is retained, not independently re-proved.
PARI candidate extraction supplies only rational unit candidates
and integer selectors.  No complete unit-basis assertion, provisional
class-group conclusion, GRH, or \texttt{bnfcertify} is used.
